\section{Propagation of the polygon laws}\label{sec:propagation}

We use the presented group $\widehat G$ and the lifting calculus of
Section~\ref{sec:lifting}.  The input consists of the true tables for a
\emph{finite} collection of geometric configurations, together with their
common translates.  Our goal is the actual central law for every finite
family of polygon tests.  We first propagate the coordinate laws by an
induction on the range of their integer labels.  We then use these
coordinate laws to compare sloping decisions on arbitrarily small
rectangular cells.  In both arguments, comparisons take place in
$\widehat G$ itself.

We retain the notation $\rho_{P,I}:\Alt(I)^\ast\to\widehat G$ for a
canonical perfect copy on a region $P$ and alphabet $I$.  When $P$ denotes
a formal sector, its rectangular coordinate will always be specified.
Saying that this copy is \emph{controlled in} a rectangle $J$ has the
meaning of Lemma~\ref{lem:control}: permutations conditional on $J$ act on
the copy as their unconditional counterparts do.  In particular, this is
an assertion about conjugations in $\widehat G$, not just a containment
of projected supports.

\subsection{The fixed geometric data}

Write $T=\RR^2/\ZZ^2$ for the ordinary coordinate torus used to discuss
local geometry.  For $z=p+q\tau\in\OO$, put
\[
 \nu(z)=q=\frac{z-\bar z}{\sqrt5}.
\]
The function $\nu$ is unchanged by adding an integer, so it also labels
an endpoint in $\OO/\ZZ$.  Let $c,C>0$ be the separation and mesh
constants of Lemma~\ref{lem:arithmetic}.  We use the choice
\begin{equation}\label{eq:prop-slope-choice}
 \lambda c>1000(C+1),\qquad |\bar\lambda|<10^{-3}.
\end{equation}
All constants below may depend on this fixed $\lambda$.

Here is the finite geometric information required from the initial
tables.  A \emph{chart} consists of a small rectangle with sides at
$\OO$ levels, one allowable line crossing it, and its two sign regions
inside the rectangle.  The sign regions are false outside the rectangle.
Include the rectangle and its coordinate quadrants in its table.  Choose
an inner rectangle strictly inside the chart, leaving a positive
ordinary margin.  We use finitely many such sizes, with nested inner
rectangles where a comparison needs a smaller margin.  Denote by
$\delta>0$ a lower bound on the margins used for comparisons in the
coordinate induction.

For a line of slope $\lambda$, its allowable tangent translations are
\[
 \{(u,\lambda u):u\in\OO\}.
\]
Choose a $\ZZ$-basis $(\eta,\lambda\eta)$,
$(\tau\eta,\lambda\tau\eta)$ with $\eta=\tau^{-b}$ so small in
ordinary norm that each step is less than $\delta/100$.  Include the
joint tables for the two charts before and after each of these steps
and its negative, and a fixed inner rectangle of their overlap.
The same choices are made for the other sloping direction, with the
coordinates interchanged.  On the overlap, corresponding sign
decisions agree.  All these assertions are literal identities in the
finite pointwise tables.

We also need finitely many charts for several concurrent lines.
Lemma~\ref{lem:arithmetic} supplies a positive integer $D$ such that an
intersection of two nonparallel allowable lines belongs to
$D^{-1}\OO^2$.  Choose representatives of the finite quotient
$D^{-1}\OO^2/\OO^2$ in a small neighborhood of $0$.  For each
representative $r$ and each subset of directions whose lines through
$r$ have allowable levels, include the true sector table of those
lines in an $\OO$-sided rectangle containing $r$ with a fixed margin.
The rectangle is chosen about $0$; its sides are not obtained by adding
$\OO$ lengths to a possibly nonintegral coordinate of $r$.
Include every coordinate clipping decision that occurs in this table.
For each sloping line in it, choose an $\OO^2$ point on that line close
to $r$, and include the comparison with a chart based at that point.
Such a point exists by density of the tangent group.  There are finitely
many choices here.  Single-line models use a nearby $\OO^2$ point on
the line, and the model with no line is constant.

Two further finite additions handle chart boundaries.  First, choose a
finite set of tangent translations that brings an inner chart close to
every point of the portion of a line meeting a fixed small enlargement
of any of the finitely many chart boxes, using its chosen $\OO^2$
anchor for a concurrent template.  The enlargement also covers the
tangential projections of nearby points at a box edge or corner.
One sufficiently short nonzero tangent translation and finitely many
of its multiples suffice.  Include the corresponding comparisons with
the original chart; their overlap rectangles are intersected with the
original box when necessary.  Second, include a fixed sufficiently
fine rectangular grid in each primitive chart.  For any fixed positive
distance from its line, this grid can be chosen so that each grid cell
meeting the part at that distance lies entirely on its correct side.
These are genuine individual tables for the chart and the grid.
This last choice will be used only away from the line; close to the
line we will use a quadrant with a vertex on the line.

Every object just described is a fixed polygon and is a Boolean
expression in finitely many translated primitive tests.  Thus the
finite-table construction and Lemma~\ref{lem:translated-copies} provide
all these tables, first on the required bounded alphabets enlarged by
the fixed spare banks, and then on every allowed subalphabet.  Bounded
integer changes of planar lift are included in the same finite list.
We will also impose one finite coordinate window at the starting size
$n_0$.  The order in which $n_0$ and the remaining parameters are fixed
is specified at the end of the proof; none of the tables depends on a
later induction level or on a word being tested.

\subsection{An overlap bridge using only the old windows}

For an interval $W$ of consecutive integers, let $\mathcal D(W)$ be
the circular partition with endpoints $i\tau\bmod1$, $i\in W$.
It is generated by the translated primitive interval tests whose two
endpoints have consecutive labels in $W$.  To see that these tests
distinguish its cells, join two purportedly indistinguishable cells by
an oriented arc.  For each consecutive pair of labels, either both
endpoints or neither must have been crossed.  Connectedness of the
integer interval $W$ says that either every endpoint or no endpoint
was crossed.  In either case the two circular cells coincide.

Write $\mathcal R(n)$ for the actual central law for the coordinate
tests with endpoint labels in arbitrary windows of at most $n$
consecutive integers, one window in each coordinate.  The statement
includes its canonical partitions, their unions of cells, and all
subalphabets.  Common translations allow the windows to start
anywhere.  For $W=[a,b]\cap\ZZ$, write
\[
 W^{+B}=[a-B,b+B]\cap\ZZ.
\]
A \emph{label margin at most $B$} means that the auxiliary endpoint
labels needed for a comparison lie in $W^{+B}$.  Thus adjoining those
labels enlarges the source range to at most $|W|+2B$ labels.  These
fixed margins will be absorbed in the strict window inequalities below.

The following lemma gives the precise transfer needed when a chart
has large labels.  It is useful even if the original interval uses
an extreme label of its source window.

\begin{lemma}[Transfer between old windows]\label{lem:window-bridge}
Assume $\mathcal R(n)$.  Fix $A,K,\delta>0$ and an integer $B\ge1$.
Suppose a perfect copy is controlled in a rectangle $R=I_x\times I_y$
whose side lengths are at most $A/n$.  Each $I_\alpha$ is a union of
cells of $\mathcal D(W_\alpha)$, where
\[
 \lfloor n/4\rfloor\le |W_\alpha|\le\lfloor0.9n\rfloor,
\]
and the source control is available with a label margin at most $B$.
Let $J=J_x\times J_y$ be a chart rectangle whose coordinate endpoint
labels each lie in an interval of at most $B$ integers.  Suppose that
these labels are at distance at most $Kn+B$ from a label of the
corresponding source window, and that $J$ contains the closed
$\delta$-neighborhood of $R$ in the chosen circular coordinate charts.
For all sufficiently large $n$, depending only on $A,K,\delta,B$,
the copy is controlled in $J$.  The proof uses only $\mathcal R(n)$.
\end{lemma}

\begin{proof}
It suffices to transfer one coordinate at a time.  Put
\[
 L=\lfloor n/8\rfloor,\qquad S=\lfloor L/2\rfloor.
\]
For $n\ge\max(160,40B)$, we have $L\ge n/9$ and $S\ge n/20$.
Choose an $L$-window $V_0\subset W$ and enclose $I$ by the cells of
$\mathcal D(V_0)$ which meet it; call their union $E_0$.
Each end of $I$ is enlarged by at most $C/L\le9C/n$.
The containment $I\subset E_0$ is compared within $W^{+B}$, whose
cardinality is at most $\lfloor0.9n\rfloor+2B<n$.  Consequently it gives
control by $E_0$ using the old law.  This first comparison discards the original
endpoint labels: no later comparison needs to retain an extreme
endpoint of $W$.

Choose an $L$-window $V_T$ containing the endpoint labels of the target
interval and their prescribed margin.  Its starting label can be
reached from that of $V_0$ in steps of at most $S$, with
\[
 T\le \lceil20(K+2)\rceil
\]
after increasing $n$ to absorb the fixed $B$ terms.  Write the
successive windows as $V_0,\ldots,V_T$.  The integer interval spanned
by consecutive windows has length at most $L+S$; even after adding
the fixed margins its length is less than $n$.  Having obtained
$E_{j-1}$, enclose it by cells from $\mathcal D(V_j)$ to obtain $E_j$.
Both enclosures occur in the old law on the interval spanned by
$V_{j-1}\cup V_j$.  Thus they have their common canonical
representations there, and $E_{j-1}\subset E_j$ transfers control.
Each endpoint grows by at most $9C/n$ at this step.

It follows that $E_T$ enlarges $I$ at each end by at most
\begin{equation}\label{eq:bridge-error}
 \frac{9C(T+1)}{n}.
\end{equation}
Choose $n$ so large that this quantity is less than $\delta$.
Then $E_T\subset J_\alpha$.  The terminal comparison is in the old
law on $V_T$ with its margin.  Transitivity of control proves the
claim for that coordinate.  A circular interval crossing the chosen
seam is simply a union of cells; the same argument is made in a
coordinate lift and then projected to the circle.  The estimate and
the available old laws are unchanged.

Apply this construction to the two coordinates.  First change one
coordinate window while keeping the other fixed, and then change the
other.  Every comparison uses an old window in each coordinate.
Finally intersect the two coordinate controls using
Lemma~\ref{lem:control}.  This proves control in $J$.
\end{proof}

The important bound in Lemma~\ref{lem:window-bridge} is the number
$T=O(K)$ of transfers, not the size of the endpoint labels.  We shall
apply the lemma separately to every overlap rectangle encountered in
a long chain of charts.  In each application $I$ is the \emph{original}
interval.  The enclosure produced for one overlap is never used as
the source interval for the next overlap.

\subsection{Growing the coordinate laws}

\begin{lemma}[Nested cuts]\label{lem:nested-cuts}
Let $E$ be a perfect group, and suppose $f,l_i,h_i:E\to H$ are
homomorphisms satisfying $f(s)=l_i(s)h_i(s)$ and
$[l_i(E),h_i(E)]=1$.  Order the indices by increasing cuts, and
suppose additionally that $[l_i(E),h_j(E)]=1$ whenever $i<j$.
Then the successive differences $l_{i-1}(s)^{-1}l_i(s)$ are commuting
homomorphisms, with the evident initial and final factors, and give
a split representation of the consecutive intervals.
\end{lemma}

\begin{proof}
For $i<j$, conjugation by $l_j(t)$ on $l_i(E)$ is conjugation by
$f(t)$, because $h_j(E)$ commutes with $l_i(E)$.  The latter
conjugation is conjugation by $l_i(t)$, because $h_i(E)$ commutes
with $l_i(E)$.  Therefore $l_i(t)^{-1}l_j(t)$ centralizes $l_i(E)$.
Writing $l_j=l_i(l_i^{-1}l_j)$ now shows directly from the
homomorphism identities that $s\mapsto l_i(s)^{-1}l_j(s)$ is a
homomorphism.  For successive indices its image centralizes every
earlier prefix, since both relevant prefixes act there by the same
conjugation $\operatorname{Ad}f(t)$.  It consequently centralizes
every earlier difference.  Set the first prefix equal to $1$ and
the last equal to $f$ if they were not already listed.  Multiplying
the differences telescopes to $f$, as required.
\end{proof}

\begin{proposition}[All rectangular laws]\label{prop:rectangular-law}
For a sufficiently large fixed initial window $n_0$, its true table
and the finite geometric tables described above imply
$\mathcal R(n)$ for every $n$.  In particular every finite family
of aligned rectangular tests has its actual central law in
$\widehat G$.
\end{proposition}

\begin{proof}
We prove the implication
\[
 \mathcal R(n)\ \Longrightarrow\
 \mathcal R(\lfloor6n/5\rfloor)
\]
for every sufficiently large $n$; iteration from $n_0$ gives
unbounded window lengths.  Put $n'=\lfloor6n/5\rfloor$.
In a window of length $n'$ on the $x$ axis take the two end
subwindows $W_-$ and $W_+$, each of length $\lfloor0.9n\rfloor$.
Their intersection has length at least $0.59n$ for large $n$.
Its circular cells form an anchor partition of mesh at most $2C/n$.

Fix an anchor cell $U$.  Each of the two old laws decomposes its
canonical copy into prefixes and suffixes at that subwindow's cuts.
Both give the same copy on $U$, since the common anchor partition
already has its old law.  Contributions on different anchor cells
commute, including contributions obtained from different subwindows:
move one small alphabet to unused letters conditionally on its
anchor cell, as in Lemma~\ref{lem:control}.  For a cut $a$ in $U$
write its prefix and suffix representations as $l_a,h_a$; these
have common star source $\Alt(I)^\ast$.  If the cut occurs in both
subwindows it belongs to the anchor partition, and the two
representations already agree.  Lemma~\ref{lem:nested-cuts} shows
that it remains to prove
\begin{equation}\label{eq:prefix-suffix}
 [l_a(\Alt(I)^\ast),h_b(\Alt(I')^\ast)]=1\qquad(a<b)
\end{equation}
for arbitrary small alphabets, and then to apply
Lemma~\ref{lem:small-support}.  Pairs coming from one subwindow
already commute by its old law.

Suppose, then, that $a$ and $b$ come from different subwindows.
The two embeddings of $\lambda$ have complementary roles in this
comparison: its large ordinary value separates two thin rectangles,
and its small conjugate value keeps the new transverse label in an old
window.
Use an ordinary lift of $U$ to order them.  The endpoint separation
bound gives
\begin{equation}\label{eq:gap-ab}
 b-a\ge c/n'.
\end{equation}
Split both copies by a $y$-grid from $\lfloor n/4\rfloor$
consecutive labels.  This is permitted separately for each copy
by the old rectangular law.  Contributions on different
$y$-cells commute by the conditional move on those cells.
Thus fix one lifted $y$-cell $[t,t+h]$, where $h\le5C/n$.
The two rectangles to be compared are
\[
 R_-=(U\cap\{x<a\})\times[t,t+h],\qquad
 R_+=(U\cap\{x>b\})\times[t,t+h].
\]
They have diameter $O(1/n)$.  By
\eqref{eq:prop-slope-choice} and \eqref{eq:gap-ab},
\begin{equation}\label{eq:slope-separation}
 h<\lambda(b-a).
\end{equation}
The allowable line
\begin{equation}\label{eq:separating-line}
 y-t=\lambda(x-a)
\end{equation}
therefore separates these rectangles: $R_-$ is above it and
$R_+$ is below it.

This geometric separation must now be proved at the level of
controlled copies.  Put
\[
 z_-=(a,t),\qquad z_+=(b,t+\lambda(b-a)).
\]
Both points lie in $\OO^2$ in the chosen planar lifts.
In the chart at $z_-$, the northwest quadrant lies above
\eqref{eq:separating-line}; in the chart at $z_+$, the southeast
quadrant lies below it.  For large $n$, both rectangles lie in
the inner chart boxes.  The actual quadrant tables, followed by
intersection control, show that the copy on $R_-$ is controlled
in the first positive sign and the copy on $R_+$ in the second
negative sign.  To justify using these tables we check the old
coordinate windows explicitly.  At $z_-$ the cuts $a,t$ are
already available, and adding fixed chart sides to a window of
length at most $0.9n$ still leaves length less than $n$.
At $z_+$ the new $y$-label is governed by
\begin{equation}\label{eq:label-contraction}
 \nu(\lambda d)=\bar\lambda\nu(d)+\nu(\lambda)d,
 \qquad d=b-a.
\end{equation}
Here $|\nu(d)|\le n'$ and $|d|\le2C/n$.
Consequently
$|\nu(\lambda d)|\le|\bar\lambda|n'+O(1)$.
Adding this displacement and the fixed chart margins to the
transverse window of length $n/4$ gives a window of length less
than $n$.  The $x$ coordinates use the second subwindow with a
fixed margin.  Every quadrant and box containment just invoked
is thus a consequence of $\mathcal R(n)$.

It remains to identify the positive sign in the two charts on the
copy on $R_-$.  Write the tangent displacement as an integral
combination of the fixed short tangent basis.  Since its ordinary
coordinates are $O(1/n)$ and its conjugate coordinates are $O(n)$,
the sum of the absolute values of these two coefficients is
$O(n)$.  This follows by applying the fixed inverse of the
two-embedding basis matrix.  Order the signed increments so that
the successive positions stay within one largest step of the
segment from $z_-$ to $z_+$.  Such an ordering exists: whenever a
partial tangent coordinate is below the interval joining the
endpoints, take a positive remaining increment; above the interval
take a negative one.  The required sign exists because the sum of
the remaining increments leads to the final endpoint.  Inside the
interval take either available sign.  Each crossing overshoots an
endpoint by at most one increment.

By the choice of the small basis, the successive charts therefore
have overlap rectangles containing both $R_-$ and $R_+$ with
one fixed positive margin, independent of $n$.  The labels of all
chart centers, relative to the starting labels, are bounded by
$Kn$ for a fixed $K$: there are $O(n)$ increments, each with
fixed labels.  Consecutive sign decisions and their overlap
rectangle belong to one of the fixed tables after a common
translation.

Apply Lemma~\ref{lem:window-bridge} to the original $R_-$ and
\emph{each} overlap rectangle separately.  The source windows
have lengths $0.9n$ and $n/4$ as required, and the overlap has
the fixed margin just obtained.  Thus the copy on $R_-$ is
controlled in that overlap.  The two sign decisions agree there
in the fixed table, so Lemma~\ref{lem:control} says that their
conjugations on the copy on $R_-$ agree exactly.  Composing
these equalities identifies the first positive sign's action
with the last one's action.

Although the chart chain has $O(n)$ steps, its geometric error
does not grow with that number.  The only enlargement is
\eqref{eq:bridge-error}, separately for each overlap and always
starting from $R_-$.  The chart chain composes equalities of
actions, rather than successively enlarged rectangles.
Figure~\ref{fig:propagation-bridge} summarizes the separation and
the short-window transfer used for each overlap.
The copy on $R_-$ is now controlled in the last positive sign,
whereas the copy on $R_+$ is controlled in its negative sign.
These two signs have a true split table.  The disjoint-control
part of Lemma~\ref{lem:control} proves their commutation and
hence \eqref{eq:prefix-suffix}.

Lemma~\ref{lem:nested-cuts} now splits all the new cuts of $U$
into their actual consecutive intervals.  Their perfect factor images
commute.  The central kernel of each source
$\Alt(I)^\ast\to\Alt(I)$ has central image in its own factor,
and hence in the product of all factors.  Every original input copy
is the prescribed product of these interval copies.  Quotienting their
generated group by its center therefore gives a quotient of the product
of ordinary alternating groups indexed by the actual intervals.
Doing this in every anchor cell supplies exactly the assignment-group
map for the new one-coordinate central law.  The identical
argument with $x-\lambda y$ proves the $y$ law.  Every pair of
one primitive $x$ test and one primitive $y$ test has a joint
true table at arbitrary offsets: a translation in the two
coordinates changes the offsets independently.  Apply
Lemma~\ref{lem:pair-laws} to assemble the new coordinate laws.
Their already established one-coordinate exclusions leave
exactly the products of actual interval cells, so no spurious
rectangular assignment remains.  Lemma~\ref{lem:small-support}
extends this law simultaneously to the required alphabets,
completing the induction.
\end{proof}

\begin{figure}[t]
\centering
\begin{tikzpicture}[x=1cm,y=1cm,>=Stealth]
 \fill[blue!13] (0.2,0.5) rectangle (2,1.6);
 \fill[orange!20] (3.1,0.5) rectangle (4.9,1.6);
 \draw[thick] (0.2,0.5) rectangle (4.9,1.6);
 \draw[dashed] (2,0.15)--(2,2.4);
 \draw[dashed] (3.1,0.15)--(3.1,2.4);
 \draw[very thick,red!65!black] (1.75,0.1)--(3.25,2.5);
 \fill (2,0.5) circle (1.3pt);
 \fill (3.1,2.26) circle (1.3pt);
 \node[below] at (2,0.05) {$a$};
 \node[below] at (3.1,0.05) {$b$};
 \node at (1,1.05) {$R_-$};
 \node at (4,1.05) {$R_+$};
 \node[left] at (0.2,0.5) {$t$};
 \node[left] at (0.2,1.6) {$t+h$};
 \node[above right] at (3.1,2.26) {$z_+$};
 \node[below left] at (2,0.5) {$z_-$};
 \draw[->,thick] (5.6,1.3)--(6.3,1.3);
 \node[align=center] at (8.35,2.6) {Integer-label windows\\for one chart overlap};
 \draw[<->,thick] (6.5,2.05)--(10.2,2.05);
 \node[left] at (6.5,2.05) {$\ZZ$};
 \draw[very thick,blue!60!black] (7.1,1.6)--(8.15,1.6);
 \draw[very thick,blue!60!black] (7.6,1.15)--(8.65,1.15);
 \draw[very thick,blue!60!black] (8.1,0.7)--(9.15,0.7);
 \node[right] at (8.15,1.6) {$V_0$};
 \node[right] at (8.65,1.15) {$V_1$};
 \node[right] at (9.15,0.7) {$V_2$};
 \draw[decorate,decoration={brace,amplitude=4pt,mirror}]
   (7.1,0.3)--(8.65,0.3);
 \node[below] at (7.875,0.16) {$|V_0\cup V_1|<n$};
\end{tikzpicture}
\caption{The sloping line separates a prefix from a later suffix.
Each geometric overlap receives a separate application of
Lemma~\ref{lem:window-bridge}, starting from the original rectangle.
Within that transfer, short integer windows may leave the source
window; the hull of each consecutive pair, with its fixed margins,
has length less than $n$.  The original extreme endpoint is discarded
in the first enclosure.  The drawing of the slope is schematic.}
\label{fig:propagation-bridge}
\end{figure}

\subsection{Comparing sloping decisions on controlled cells}

We have now established all rectangular laws, without assuming any
joint law for arbitrarily translated sloping predicates.  Consequently
a rectangle and any of its rectangular containers can be compared
exactly, with no restriction on labels.  This removes the quantitative
window issue from the rest of the proof.

\begin{lemma}[Comparison along one line]\label{lem:line-slide}
Let $p$ and $q$ be two chart decisions for the same oriented allowable
sloping line, either primitive charts or charts in the fixed
concurrent templates.  Suppose a point $z$ is on this line, or within
a fixed sufficiently small distance of it, and belongs to the
closures of both chart boxes.  On all sufficiently small rectangular
cells $P$ near $z$ contained in both boxes, $p(s)$ and $q(s)$ have the same
conjugation on every perfect copy controlled in $P$.  The conclusion
also holds when the copies are formal sloping sectors with actual
rectangular control.  Only the rectangular laws and the fixed
geometric tables are used.  At a clipping boundary the comparison is
made on the common interior cells; on exterior cells the decision
from that chart is zero.  Thus a replacement used on both sides must
retain the original clipping condition.
\end{lemma}

\begin{proof}
First suppose the two charts are primitive charts with $\OO^2$
vertices on their common line.  At the first chart use one of the
fixed recentering offsets to obtain an inner chart near $z$.
The first overlap is intersected with the original clipping box.
On a cell on its interior side this is a rectangular container of
$P$.  The fixed comparison table proves equality of the two
decisions on that container.  The same construction is made at
the final chart.  Exterior cells are not an assertion of equality
of the unmasked signs: the original decision is zero by its
actual box-complement law, and a replacement there is clipped
by that same box.

The two recentered $\OO^2$ vertices differ by an element of the
tangent group.  Express this difference in the fixed short basis
and order the signed increments as in the proof of
Proposition~\ref{prop:rectangular-law}.  The centers remain in a
small neighborhood of the segment joining the vertices and hence
near $z$.  More precisely, choose a fixed $\varepsilon>0$ much
smaller than every inner chart margin.  The finite recentering
nets may be chosen to put both endpoints within $\varepsilon/10$
of the tangential projection of $z$, and each basis step may be
chosen shorter than $\varepsilon/10$.  The greedy ordering keeps
every intermediate center within $\varepsilon/2$ of that projection.
If the distance of $z$ from the line and the diameter of $P\cup\{z\}$ are
less than $\varepsilon/10$, every middle overlap therefore has a
fixed positive margin around $P$, independently of the number of
steps and all labels.  The first and last overlaps have the same
margin in their uncut directions; they are simply clipped at the
original box edges.  This supplies a uniform smallness bound,
without shrinking a cell once for each step of a long chain.
All rectangular containers now have their laws by
Proposition~\ref{prop:rectangular-law}, irrespective of their
labels.  Every consecutive fixed table therefore gives equal
actions on a controlled copy by Lemma~\ref{lem:control}.
Their composition gives the desired equality.

For a template with vertex $r+u$, where $u\in\OO^2$, use the
chosen $\OO^2$ anchor on each line of its representative at $r$,
translated by $u$.  The initial and terminal primitive anchors
lie on the same labeled line; their difference is exactly in its
tangent group.  The fixed terminal comparison identifies the
last primitive chart with the template chart.  This accounts
explicitly for vertices in $D^{-1}\OO^2$ without pretending that
the vertex itself is an allowable translation.  A single line
through an arbitrary ordinary point needs only an $\OO^2$
anchor on that line sufficiently close to the point.

At either clipped endpoint every overlap is intersected with
the appropriate inside coordinate half-planes of its box.
These are rectangular conditions with their established law.
Thus the argument applies separately to each actual rectangular
side of the clipping boundary.  It does not require the controlled
copy to split by any intermediate sloping sign.  Its sole use of
the copy is its control in the rectangular overlap.
\end{proof}

\begin{lemma}[One sloping predicate and rectangular partitions]
\label{lem:slope-rectangle}
A translated primitive sloping predicate has a joint central law
with every finite aligned rectangular partition.  At this stage
one may retain both formal sloping sides over each rectangular
cell.  Every resulting perfect copy has the actual rectangular
control of that cell, and hence control in all its rectangular
containers.
\end{lemma}

\begin{proof}
We first compare one sloping predicate $p$ with one translated
primitive coordinate predicate $v$.  Around $p$ choose the fixed
fine rectangular grid included in its initial table.  Include
its clipping edges and its planar chart seams in the grid.
The mesh is chosen smaller than the gaps between the two
boundaries of a primitive coordinate predicate, and sufficiently
small compared with the fixed chart margins and the finitely
many angles.  Translate the whole grid with $p$.

This partition has an actual joint law separately with $p$, by
the fixed table, and with $v$, by the rectangular law.  On one
of its cells $P$, form $H_P$ from the restricted copies of these
two predicates and their complements, as in
Lemma~\ref{lem:patching}.  If either predicate is constant on
$P$, its action on $H_P$ is the indicated constant action.  This
assertion follows from its \emph{individual actual} split with
$P$, together with the conditional move on $P$; it does not use
the pair law being proved.  Only the other individual split is
then needed to check a formal two-predicate relation.

If both predicates vary, $P$ is inside the clipping box of $p$,
one sloping line crosses it, and exactly one boundary line of
$v$ crosses it.  Their intersection is within a fixed multiple
of the diameter of $P$: solve the two line equations, whose
angle is one of a fixed finite list.  It lies in
$D^{-1}\OO^2$.  The grid was chosen fine enough that the
appropriate translated concurrent template has a margin box
containing $P$.  The intersection itself can lie just outside
the clipping box of $p$.  In applying the line comparison choose
instead a point $z_0$ of its sloping line in $\overline P$.
The template intersection is within $C_\lambda\operatorname{diam}P$
of $z_0$, and its chosen $\OO^2$ anchor is within a fixed arbitrarily
small distance of that intersection.  Choose the mesh once so
that this distance and the recentering errors fit the uniform
margin established in Lemma~\ref{lem:line-slide}.  The tangent
walk then stays near $P$ regardless of its number of steps.
No refinement depending on the relative label of $v$ is needed.
Replace $p$ by its corresponding line decision
in that template, using Lemma~\ref{lem:line-slide}, and replace
$v$ by the matching axis decision, using only rectangular
control.  The replacements have exactly the same actions on
$H_P$ as the original generators.

The true template has the four ordinary sectors of these two
crossing lines.  Extend it outside its margin box by any one
of those four assignments.  Therefore a word trivial on all
formal assignments of $(p,v)$ is trivial in this pointwise
template model.  Its lifted value is central in the template
subgroup.  Its rectangular margin box controls every generator
of $H_P$, since it contains $P$ and all rectangular laws are
available.  Take the template law on the enlarged alphabets
required by Lemma~\ref{lem:patching}.  The template-relator
conclusion of that lemma now makes the replacement word act
trivially on $H_P$.

Thus every formal two-predicate relator acts trivially on each
$H_P$.  Lemma~\ref{lem:patching} and then
Lemma~\ref{lem:small-support} prove the pair law.  Apply this
to the coordinate predicates generating any desired rectangular
partition, and apply Lemma~\ref{lem:pair-laws}.  Retain all
exclusions already supplied by the rectangular law.  The
resulting sectors have an actual rectangular cell coordinate,
although their sloping side may still be formal.  Refining by
another rectangle and using the same pair laws proves that
each sector is controlled in every rectangular container of
its cell.  No compatibility of two different sloping predicates
has been used.
\end{proof}

For a finite family $p_1,\ldots,p_k$ and a common rectangular
partition $\mathcal P$, the lemma supplies the separate laws
needed to form the cell groups of Lemma~\ref{lem:patching}:
$H_P$ is generated by all the $p_j$-side factors over $P$.
That lemma makes every $H_P$ invariant under each $p_j(s)$,
so equalities of their actions can be composed before a joint
sloping law is known.  Every generator of $H_P$ is controlled
in $P$ and its rectangular containers, including a generator
assigned a geometrically impossible sloping side.  These are
the groups on which we will verify the actual polygon laws.

\subsection{Actual sectors and clipping boundaries}

The remaining task is to remove those formal extra sides and
prove the joint law for arbitrarily many sloping predicates.
The small rectangular cells are now allowed to depend on the
finite family or word being considered.  The available tables
and the presented group remain fixed.

\begin{lemma}[An inactive line gives a constant action]
\label{lem:inactive-line}
Let $p$ be a translated primitive sloping predicate with clipping
box $B$, and let $z$ be an ordinary point not on its sloping
line $L$.  For sufficiently small rectangular cells $P$ near
$z$ respecting the edges of $B$, the action of $p(s)$ on any
perfect copy controlled in $P$ is its actual constant-side
action inside $B$, and is trivial outside $B$.  The assertion
holds for the formal sector copies of
Lemma~\ref{lem:slope-rectangle}.
\end{lemma}

\begin{proof}
Outside $B$ the actual law of $p$ with its own clipping box
gives the assertion, by exterior rectangular control.
Consider the interior side, allowing $z$ itself to be on an
edge or corner of $B$.  We give the argument for
$L=\{y-\lambda x=c_0\}$; interchange coordinates for the
other slope.  Put $d=z_y-\lambda z_x-c_0\ne0$.

Suppose first that $d>0$.  The point on $L$ with coordinates
\[
 (z_x+d/(2\lambda),\ z_y-d/2)
\]
has $z$ strictly northwest of it.  By density choose
$q_x\in\OO$ within $d/(4\lambda)$ of its first coordinate,
and put $q_y=\lambda q_x+c_0\in\OO$.  Then
\begin{equation}\label{eq:strict-quadrant}
 q_x-z_x>d/(4\lambda),\qquad z_y-q_y>d/4.
\end{equation}
The northwest quadrant based at $q=(q_x,q_y)$ lies wholly on
the positive side of $L$.  For $d<0$, the same formula with
$|d|/(4\lambda)$ as the allowed perturbation puts $z$ strictly
southeast of $q$; that quadrant lies wholly on the negative
side.  Thus the coordinate inequalities have positive room
even when the inactive line is arbitrarily close to $z$.

If $|d|$ is below a fixed small threshold, choose that threshold
so that $q$ and $z$ fit in the inner margins of a line chart.
Lemma~\ref{lem:line-slide} compares the original predicate with
that chart near $z$, clipping the first overlap to $B$.
By \eqref{eq:strict-quadrant}, all sufficiently small cells on
the interior side of $B$ are controlled in the indicated
quadrant and in the chart box.  The genuine line--quadrant
table says that its sign decision there is respectively the
constant permutation or the identity.  Intersection and
transitivity of control give the claimed action on the
controlled copy.  No sign attached to a formal sector was
used to infer its location.

For $|d|$ above the fixed threshold, use the finite rectangular
grid included with the primitive chart.  Choose its mesh small
enough that every cell meeting the part with this lower bound
on $|d|$ lies strictly on the same side of the line.  This is
possible because the defining linear form has fixed norm.
The individual actual table of $p$ with the grid gives its
constant action there.  Rectangular control transfers the
assertion to any sufficiently small $P$ near $z$.  Cells
meeting a grid boundary may be further split by that fixed
grid; this changes neither the conclusion nor the initial
data.

If a clipping edge is active at $z$, treat interior cells by
the quadrant or grid just described and exterior cells by
the first paragraph.  At a corner there are four rectangular
sides and the same procedure applies to each.  Therefore the
local action is exactly the Boolean combination of the
constant sloping sign with the actual clipping decisions on
all sides of the boundary.
\end{proof}

\begin{theorem}[The polygon law]\label{thm:polygon-law}
There is one finite choice of true initial relations in
$\widehat G$, for sufficiently large fixed finite $m$, such
that every finite family of aligned polygon tests has its
actual central law.  The assertion holds simultaneously on
subalphabets.  Its canonical representations
\[
 \rho_{U,I}:\Alt(I)^\ast\longrightarrow\widehat G
\]
are compatible with restriction of the alphabet, split
exactly over finite disjoint polygon refinements, and are
covariant under common translations on proper supports.
\end{theorem}

\begin{proof}
First take a finite family of translated primitive predicates,
including coordinate predicates if needed.  By
Lemma~\ref{lem:small-support}, we may fix an alphabet of its
prescribed bounded size and a word $w$ entirely on that alphabet
whose pointwise permutation is the identity on every actual
assignment.  It suffices to prove that $w$ acts trivially on
the subgroups needed to test centrality.  Include those
additional small-support copies and their spare alphabets
throughout; the geometric family is still finite.

Fix $z\in T$.  Call a supporting boundary line \emph{active}
at $z$ if it contains $z$; include the coordinate lines which
form clipping edges.  Inactive supporting lines have positive
distance from $z$, and the family is finite.  In a sufficiently
small ordinary neighborhood of $z$ the only varying signs are
therefore the active ones.  There is at most one distinct
active line of each of the four directions.

Choose the corresponding translated concurrent template.
If at least two directions are active, $z\in D^{-1}\OO^2$
and the template is provided by its coset representative and
direction subset.  With one active direction use a nearby
$\OO^2$ anchor on that line, and with none use a constant
model.  Each original predicate becomes, in this model, the
Boolean combination of its active sign and clipping decisions,
with the other decisions replaced by their actual local
constants.  The allowed assignments are precisely the
ordinary open sectors of the active lines, with these
combinations evaluated on them.  Each such sector occurs
arbitrarily near $z$ in the original Boolean space.  Hence
the word $w$ is the identity on every assignment of this
template model.  Extend the model outside its margin box
by any one allowed sector, so that it remains a finite
pointwise test model on the whole square.

We verify the equality of its actions with the original
actions on $H_P$ for every sufficiently small cell near $z$.
All cells will be cut by every original clipping edge and
by any additional finite rectangular subdivisions needed
in these comparisons.
\begin{enumerate}[label=(\roman*)]
\item Coordinate decisions compare with the template by
the established rectangular law.  An inactive coordinate
decision is its actual constant on the relevant cells.
\item An active sloping line inside its clipping box
compares with the same oriented line in the template by
Lemma~\ref{lem:line-slide}.  If $z$ lies on a clipping edge,
make that comparison only on the interior rectangular
side.  The template includes that active axis decision,
and its true table identifies the resulting clipped sign
there.  On exterior cells the original predicate is
trivial by control in the exterior of its own box.
This treats all sides of a corner as well.
\item An inactive sloping line is replaced by its actual
constant-side decision by Lemma~\ref{lem:inactive-line}.
If some of its clipping edges are active, keep those
coordinate decisions in the replacement.  The interior
and exterior arguments in that lemma establish precisely
this clipped Boolean combination, even for a formal
wrong-side copy over $P$.
\end{enumerate}
Every equality in this list follows from equal actions on
a known rectangular container of $P$ and the conditional
move of Lemma~\ref{lem:control}.  Thus it holds on every
generator of $H_P$, not just in projection.  The original
generators preserve $H_P$ by Lemma~\ref{lem:patching}.  Their actions
can therefore be composed, and the action of $w$ on $H_P$
equals the action of its template replacement.

The replacement is central in the true template table.
Its margin box is a rectangular container of $P$, so it
controls every generator of $H_P$.  Apply the table on the
enlarged alphabets required by Lemma~\ref{lem:patching}.
The template-relator conclusion of that lemma shows that
the replacement, and hence $w$, centralizes $H_P$, including
its possibly formal sloping factors.

We finally choose one finite partition on which these
local checks all apply.  For every $z$ the preceding
construction provides an ordinary neighborhood and an
upper bound on the permitted cell diameter.  A finite
subcover exists by compactness of $T$.  A sufficiently
fine $\OO$ rectangular grid, refined by the finitely many
clipping edges and auxiliary rectangular cuts of that
subcover, has every cell in one of those neighborhoods
and below its required diameter.  This follows, for
example, by applying the Lebesgue number property to
the finite subcover and then using the coordinate mesh
bound.  Make the separate splits of every original
predicate with this partition by
Lemma~\ref{lem:slope-rectangle}.  We have proved that $w$
centralizes every resulting $H_P$.  The group generated by these
subgroups contains the original test subgroup, so
Lemma~\ref{lem:patching} proves its actual central law.
Lemma~\ref{lem:small-support} gives the assertion on all
the required alphabets simultaneously.

Apply this established law to a larger finite family containing
translated primitive coordinate tests that generate the partition just
used.  Apply Lemma~\ref{lem:law-enlargement} first to the coordinate
subfamily: its original partition copies agree with the corresponding
coordinate factors of the enlarged actual law.  We may therefore include
those same copies among the inputs of that law.  For each old formal law,
the enlarged family now contains exactly its specified input copies,
including the constant copy.  Its subgroup $H_0$ is thus contained in
the enlarged defining-copy subgroup $K$.
Every new actual assignment restricts to an allowed old assignment,
since the old law retained all actual assignments.
Lemma~\ref{lem:law-enlargement} now identifies each \emph{original}
formal factor with the product of its actual refinements and makes it
trivial when it has no actual extension.

Every polygon
is a Boolean combination of finitely many translated
primitive tests.  Apply the law just proved to a common
finite family realizing the polygons.  The canonical
perfect-cover construction supplies their representations.
To compare two Boolean expressions, pass to their common finite
family and apply Lemma~\ref{lem:law-enlargement}; the resulting
representations agree there.  Within this lawful family, uniqueness
of lifts from a perfect source gives exact
multiplicative splitting over a finite disjoint
refinement.  Finally a common allowable translation on
a proper alphabet is implemented by the balanced
translations of Lemma~\ref{lem:translated-copies}; apply
uniqueness to the translated test family.  This proves
the claimed translation covariance and alphabet
compatibility.
\end{proof}

\subsection{Why one finite presentation suffices}

For clarity, we finish by giving the order of the choices used in
this section.  First fix the arithmetic constants and choose
$\lambda$ satisfying \eqref{eq:prop-slope-choice}.  Next choose the
finite concurrent representatives, chart boxes and margins,
quadrant tables, and the fixed grids used away from a line.
Choose the short tangent bases, finite recentering offsets, and
comparison tables inside those margins.  These choices determine
a finite label bound $B$ for every fixed relative configuration,
and a finite constant $K$ for the chart walks in the rectangular
induction.  They also determine all bounded support requirements.
Choose the finite track number $m$ beyond these requirements,
the spare banks of the lifting calculus, and the thresholds
required elsewhere in the argument.

Only now choose $n_0$.  It is large enough for every strict window
inequality, for the rectangles and the separating segment to fit
in the prescribed chart margins, and for
\[
 \frac{9C(\lceil20(K+2)\rceil+1)}{n_0}<\delta.
\]
Enlarge it also for the fixed margins in
\eqref{eq:label-contraction}.  Finally impose the finite true
tables for the coordinate window of length $n_0$ and for the
finite geometric menu, with their specified representative
words and bounded alphabet enlargements.  After common
translations these are exactly the tables used above.

Later induction levels add no initial relation.  Nor does an
arbitrarily close inactive line require a new small geometric
template: its strict quadrant is a translate of a fixed
line--quadrant table, and only the independently controlled
rectangular cell becomes smaller.  Likewise arbitrarily long
words require more cells and more successive comparisons, but
only finitely many tracks at any one comparison.  Thus
Theorem~\ref{thm:polygon-law} is a statement inside a single
finitely presented group $\widehat G$.
